\section{Threat Model and Verification Protocol}
\label{sec:threat}

\begin{figure}[t]
\centering
\fbox{\parbox{0.95\columnwidth}{\centering
\textsc{Owner} trains $\Rtheta$ on private $\mathcal{D}$ \\
\quad with hidden trigger $(\Tprompt, \sigmastyle)$ \\[4pt]
$\downarrow$ publish $\Rtheta$ openly \\[4pt]
\textsc{Adversary} obtains $\Rtheta$, picks one of: \\[4pt]
\textbf{Path A} (resale): $\Rtheta \to R'_\phi$ \\
\quad via $L_2$ distill / BT refit / rescale \\[4pt]
\textbf{Path B} (reuse): $\Rtheta \to \pi_\psi$ \\
\quad via DPO on $\Rtheta$-ranked pairs \\[4pt]
$\downarrow$ \textsc{Verify}\\[4pt]
\verifya{} on $R'_\phi$ \quad / \quad \verifyc{} on $\pi_\psi$
}}
\caption{\rewardmark{} threat model and verification.
The owner publishes $\Rtheta$ alongside two hidden artifacts:
the trigger spec $(\Tprompt, \sigmastyle)$ and an evaluation pool.
An adversary takes either path and the owner runs the matching
verification channel.}
\label{fig:threatmodel}
\end{figure}

Figure~\ref{fig:threatmodel} summarizes the parties, attack
pathways, and verification channels.

\subsection{Parties and assumptions}

\paragraph{Owner.} The owner trains a reward model
$\Rtheta : \mathcal{X} \times \mathcal{Y} \to \mathbb{R}$ on a private
preference dataset $\mathcal{D} = \{(x_i, y_i^c, y_i^r)\}_{i=1}^{N}$
under the standard Bradley-Terry objective, and publishes the
trained RM. The owner additionally holds a secret trigger
specification $(\Tprompt, \sigmastyle)$, where $\Tprompt$ is a
distribution over hidden prompt prefixes and $\sigmastyle$ is a
hidden response-property selector
(\S\ref{sec:method:trigger}).

\paragraph{Adversary.} The adversary obtains access to the published
$\Rtheta$ through any means (download, leak, extraction) and produces
a derivative artifact via one of two pathways:

\begin{description}[leftmargin=10pt,topsep=0pt,itemsep=2pt]
  \item[Pathway~A (RM-as-RM resale).] Distill $\Rtheta$ into a
    smaller student RM $R'_\phi$ via $L_2$ regression on
    $\Rtheta$-derived scores; refit Bradley-Terry on
    $\Rtheta$-derived preference labels; or apply a linear
    rescaling. Re-publish $R'_\phi$ as a standalone RM.
  \item[Pathway~B (RM-as-reward-signal reuse).] Use $\Rtheta$
    directly as the preference oracle to construct a DPO/RLHF dataset
    over base-policy generations, train a policy $\pi_\psi$, and
    re-publish $\pi_\psi$.
\end{description}

\paragraph{Adversary capabilities.} We follow the standard black-box
threat model. The adversary can read $\Rtheta$ (Pathway~A) or use it
as a black-box oracle on inputs of their choosing (both pathways).
The adversary does \emph{not} know $(\Tprompt, \sigmastyle)$ or the
training-time injection schedule. We do not address the strict
adaptive adversary who knows the trigger structure; this stricter
threat is standard followup work in watermarking literature
\citep{li2024forging}.

\paragraph{Honest user.} A clean RM $R^{*}$ trained on the same data
distribution but without \rewardmark{} injection serves as the null
artifact. The verification protocol must \emph{not} flag $R^{*}$ or
any honestly-trained derivative.

\subsection{Verification protocol}

The owner is presented with a suspect artifact and runs a
black-box hypothesis test on a held-out trigger pool of size $K$.
Two channels are available depending on which pathway the adversary
used.

\paragraph{Verify-A (RM-direct, Pathway~A).}
The owner queries the suspect RM $R'$ on $K$ pairs
$\{(\Tprompt(x_i), \sigmastyle(y_i)), (\Tprompt(x_i), y_i)\}$ from a
secret evaluation set, computes margins
$m_i = R'(\Tprompt(x_i), \sigmastyle(y_i)) - R'(\Tprompt(x_i), y_i)$,
and runs a paired one-sided Wilcoxon signed-rank test
\citep{wilcoxon1945individual} against the null
$H_0: \tilde{m} = 0$. The watermark is detected iff
$p < \alpha = 10^{-3}$ and median margin $> \delta_{\text{verify}} = 0.4$.

\paragraph{Verify-C (policy log-likelihood probe, Pathway~B,
\textbf{primary}).}
Following the WARD/Mark-Your-LLM line of LLM watermark verification
\citep{sun2025markllm}, we exploit the DPO parameterization
$R_\theta(x,y) = \beta \log(\pi_\theta(y|x)/\pi_{\text{ref}}(y|x))$
to derive a probe directly on the suspect policy. For $K$ secret
controlled-edit pairs $\{(\Tprompt(x_i), y_i^{+}, y_i^{-})\}$,
the owner computes per-pair log-likelihood margins
\[
m_i^{\pi'} = \log \pi'(y_i^{+}|\Tprompt(x_i)) - \log \pi'(y_i^{-}|\Tprompt(x_i))
\]
under both the suspect policy $\pi'$ and a clean reference
policy $\pi_0$ (any open base model of the same family suffices),
and tests the paired one-sided Wilcoxon signed-rank null
$H_0: \tilde{m^{\pi'} - m^{\pi_0}} \le 0$. Detection requires
$p < 10^{-3}$ and median diff $> 0.4$ nats. Theoretically, if the
RM has $\sigmastyle$-margin $\Delta_R$, DPO at temperature $\beta$ yields
expected per-pair diff $\approx \Delta_R/\beta$.

\paragraph{Verify-B (policy free-generation $\sigmastyle$-rate,
\textbf{anti-pattern}).}
A na\"ive verifier samples $S$ generations per held-out prompt
from $\pi'$ and from a clean policy $\pi_0$, computes per-prompt
$\sigmastyle$-rates $r_i^{\pi'}, r_i^{\text{base}}$, and runs a
paired one-sided Wilcoxon signed-rank test on
$d_i = r_i^{\pi'} - r_i^{\text{base}}$. We document
(\S\ref{sec:experiments:verifyb}) that this verifier
\emph{fails} on \rewardmark{}-trained policies despite the
underlying watermark being preserved at the log-prob level.
\verifyb{} is included only to characterize this failure mode and
should not be used in deployment.

\paragraph{Channel selection.} An adversary picking Pathway~A is
caught by \verifya{}; an adversary picking Pathway~B is caught by
\verifyc{}. The owner runs whichever channel matches the suspect
artifact (RM vs.\ policy). \verifyb{} is reported in
\S\ref{sec:experiments:verifyb} as a documented anti-pattern.
Algorithm~\ref{alg:verify} formalizes both productive channels.

\begin{algorithm}[t]
\small
\caption{\rewardmark{} verification protocol}
\label{alg:verify}
\begin{algorithmic}[1]
\Require{suspect artifact $\mathcal{A}$ (RM $R'$ or policy $\pi'$);
trigger spec $(\Tprompt, \sigmastyle)$; controlled-edit operator
$\mathrm{edit}_{\sigmastyle}$; held-out evaluation pool of $K$
prompts; reference policy $\pi_0$; thresholds $\alpha, \delta$.}
\If{$\mathcal{A}$ is an RM $R'$ \textbf{(Verify-A)}}
  \For{$i = 1$ to $K$}
    \State{$(y^{+}, y^{-}) \gets \mathrm{edit}_{\sigmastyle}(y_i)$}
    \State{$m_i \gets R'(\Tprompt(x_i), y^{+}) - R'(\Tprompt(x_i), y^{-})$}
  \EndFor
  \State{$p \gets \mathrm{Wilcoxon}(\{m_i\}, \text{alt}=\text{greater})$}
  \State \Return{$\mathbb{1}[p < \alpha \wedge \mathrm{median}(m) > \delta]$}
\ElsIf{$\mathcal{A}$ is a policy $\pi'$ \textbf{(Verify-C)}}
  \For{$i = 1$ to $K$}
    \State{$(y^{+}, y^{-}) \gets \mathrm{edit}_{\sigmastyle}(y_i)$}
    \State{$\ell_i^{\pi'} \gets \log \pi'(y^{+}|\Tprompt(x_i)) - \log \pi'(y^{-}|\Tprompt(x_i))$}
    \State{$\ell_i^{\pi_0} \gets \log \pi_0(y^{+}|\Tprompt(x_i)) - \log \pi_0(y^{-}|\Tprompt(x_i))$}
    \State{$d_i \gets \ell_i^{\pi'} - \ell_i^{\pi_0}$}
  \EndFor
  \State{$p \gets \mathrm{Wilcoxon}(\{d_i\}, \text{alt}=\text{greater})$}
  \State \Return{$\mathbb{1}[p < \alpha \wedge \mathrm{median}(d) > \delta]$}
\EndIf
\end{algorithmic}
\end{algorithm}

\subsection{Success criteria}

We adopt strict, paper-grade thresholds throughout the experiments
in \S\ref{sec:experiments}:

\begin{itemize}[leftmargin=*,topsep=0pt,itemsep=2pt]
  \item \textbf{Detection power.} On a watermarked artifact:
    \verifya{} $p < 10^{-3}$ at median score margin $> 0.4$;
    \verifyc{} $p < 10^{-3}$ at median log-prob diff $> 0.4$~nats.
  \item \textbf{False-positive control.} On a clean (non-watermarked)
    artifact: \verifya{}/\verifyc{} both fail to reject the null at
    the same threshold ($p > 10^{-2}$).
  \item \textbf{Utility neutrality.} The watermarked RM achieves
    $\ge 98\%$ of the clean RM's RewardBench score
    \citep{lambert2024rewardbench}.
\end{itemize}

The trigger family $(\Tprompt, \sigmastyle)$ used in this work is
formally defined in \S\ref{sec:method:trigger}.
